\section*{Overview}

Closest pair is the standard reminder that geometry problems often become about ordering, not formulas. The naive
\(O(n^2)\) scan compares every pair of points. The divide-and-conquer solution wins by proving that after solving left
and right halves, only a tiny number of cross-half candidates still matter.

\section*{When to Use It}

Use it when:

\begin{itemize}
  \item the task asks for the nearest pair under Euclidean distance,
  \item the points are static,
  \item \(n^2\) comparisons are too expensive.
\end{itemize}

\section*{Core Idea}

Sort points by \(x\), split at the middle, solve both halves recursively, and let \(d\) be the better of the two
distances. Any cross-half answer must lie in the vertical strip of width \(2d\) around the split line.

\section*{Key Insight}

After sorting the strip by \(y\), each point only needs to be checked against a constant number of later points. The
packing argument is the heart of the algorithm: too many points inside a small \(d \times 2d\) neighborhood would force
two of them to be closer than \(d\) already.

\section*{Worked Problem}

\paragraph{Problem.}
Given \(n\) watchtowers on a plane, find the squared distance between the two closest towers.

\paragraph{Algorithm outline.}

\begin{itemize}
  \item Sort once by \(x\).
  \item Recurse on left and right halves.
  \item Merge the halves by \(y\).
  \item Build the strip and check only the next few points in \(y\)-order.
\end{itemize}

\section*{Correctness Intuition}

The recursive answers handle pairs entirely inside one half. For cross-half pairs, both points must be within distance
\(d\) of the dividing line; otherwise their \(x\)-difference alone is already too large. Inside the strip, the packing
bound limits how many candidates each point can have.

\section*{Complexity Analysis}

With the standard merge-by-\(y\) implementation, the recurrence is
\[
T(n) = 2T(n/2) + O(n),
\]
so the total complexity is \(O(n \log n)\).

\section*{Implementation}

The code keeps points sorted by \(x\), performs the recursive divide-and-conquer, and returns the minimum squared
distance as a 64-bit integer.

\section*{Common Pitfalls}

\begin{itemize}
  \item Re-sorting the strip from scratch at every recursion level and losing the \(O(n \log n)\) bound.
  \item Using floating point even though squared integer distances are enough.
  \item Forgetting duplicates; the answer can be \(0\).
  \item Mixing up the strip width when working with squared distances.
\end{itemize}

\section*{Variants / Extensions}

\begin{itemize}
  \item Manhattan-metric variants after coordinate transforms.
  \item Dynamic closest pair, which is a much harder data-structure problem.
  \item Closest pair in higher dimensions with more expensive constants.
\end{itemize}

\section*{Practice Problems}

\begin{itemize}
  \item Static Euclidean closest pair.
  \item Detect duplicates or near-collisions in large point sets.
  \item Any geometry divide-and-conquer where only a strip of cross-boundary candidates survives.
\end{itemize}

\section*{References}

\begin{itemize}
  \item \href{https://cp-algorithms.com/geometry/nearest_points.html}{cp-algorithms: Finding the Nearest Pair of Points}.
  \item \href{https://usaco.guide/plat/geo-ext}{USACO Guide: Extra Geometry}.
  \item \href{https://codeforces.com/catalog}{Codeforces Catalog}.
\end{itemize}
